\subsection{Scalability and Practicality}
Scalability has distinct mathematical, computational, and physical meanings
in this protocol. The alternative phase branch has an ideal restricted-menu
equilibrium for every $N$, and compact GHZ evaluation reaches $N=128$
without dense state enumeration. Neither result increases the number of
simultaneous players demonstrated on hardware: the native device evidence
ends at $N=7$. The new alternative-phase batch tests every D/H deviation
but does not resolve the registered tolerance at $N=7$. Ideal per-player gains
scale as $C/N$, incentive margins decrease with $N$, and the permitted
entanglement-angle interval narrows. Practical scaling therefore requires
explicit statistical precision and implementation budgets.

The protocol assumes a shared preparation stage and joint inverse/readout.
Our centralized experiments do not demonstrate a distributed quantum market
or remove the need to trust preparation, decoding, and enforcement of a
permitted strategy menu. Market architectures motivate interpretations of
the graph families, but are not validated deployment models. Comparisons
with classical mechanisms must state which mediation and enforcement
resources are available~\cite{vanenk2002classical}.

\subsection{Limitations and Future Work}
Our primary game fixes $V=4$ and $C=3$ in a strict-dominance regime and charges
conflict cost only at the all-Hawk outcome. High mean retention is therefore
compatible with random outcomes, unequal allocations, and profitable
deviations. The partial-conflict sensitivity analysis tests this modeling
choice retrospectively; it does not establish a realistic economic payoff
model or a noisy equilibrium of the altered game.

All restricted equilibria refer to explicitly named menus. Both cooperative
phase families admit a profitable unrestricted SU(2) deviation. Historical
device measurements tested Hawk deviations only, and do not establish
resistance to Dove deviations under hardware noise. Arbitrary quantum
channels, ancillary systems, coalitions, and adaptive multi-round strategies
are outside our certification scope. Symmetric strategy search also does not
exclude better asymmetric profiles or other equilibria.

Entangler-family comparisons depend on the chosen generator, interpolation
angle, compiled circuit, routing, and noise placement. The W results concern
one specified operator; equal nominal angles do not equalize interaction
budget or entangling power. The normalized interaction convention defined
above remains an unexecuted sensitivity control. Noiseless star graph-role
asymmetry and noise-induced ring circuit-order asymmetry are distinct effects;
the star wiring experiment alone does not identify the ring's mechanism.

Depolarizing noise is a controlled model rather than a complete device
description. Readout correction and zero-noise extrapolation introduce
assumptions and uncertainty; extrapolated estimates can exceed physical
payoff bounds. Chain-matched scaling uncertainty rests on two calibration
epochs, and the original $N=6,7$ batch has not been repeated. Registered
one-parameter models did not extrapolate reliably. The new complete
alternative-phase batches are repetitions within one session, and fail
their overall five-size criterion; only $N=4,6$ pass in both. The compact simulation
results do not extend the demonstrated physical player count. The
finite-round adaptation pilot remains exploratory and does not establish
convergence of general learning rules.
